% TOP_PROBLEM: {"id": 10000012, "title": "Uniform expansion bounds for graph nets", "category_id": 3, "category": {"id": 3, "name": "graph_theory", "display_name": "Graph Theory", "description": "Problems involving graphs, networks, and their properties.", "slug": "graph-theory", "order_index": 3, "created_at": "2026-07-31T15:26:25.671Z"}, "status": "open", "problem_number": "AMR-099-0012", "set_id": 13, "statusQualification": "Defines the optimal lower bound F so that the source’s arbitrary lower-bound function f cannot be chosen artificially small."}
% FIRST_POSTED: 2026-10-01
% ENTRY_KIND: statement_only
% UnsolvedMath ID 10000012; source catalogue code AMR-099-0012
% !TeX program = lualatex
% Definition and sources only; this file is not an LLM solution attempt.
\documentclass[11pt,a4paper]{article}
\usepackage[margin=25mm]{geometry}
\usepackage{fontspec}
\setmainfont{lmroman10-regular.otf}[BoldFont=lmroman10-bold.otf,
 ItalicFont=lmroman10-italic.otf,BoldItalicFont=lmroman10-bolditalic.otf]
\usepackage{amsmath,amssymb,mathtools}
\usepackage{unicode-math}
\setmathfont{latinmodern-math.otf}
\AtBeginDocument{\let\setminus\smallsetminus}
\usepackage{xurl,hyperref}
\hypersetup{colorlinks=true,urlcolor=blue}
\setcounter{secnumdepth}{0}
\setlength{\parindent}{0pt}
\setlength{\parskip}{5pt}
\setlength{\emergencystretch}{3em}
\begin{document}
\section*{Uniform expansion bounds for graph nets}\label{10000012}
{\small UnsolvedMath ID 10000012 · AMR-099-0012 · Graph Theory · AMR Open Problem Lists}
\subsection{Definitions and mathematical statement}
For an infinite connected graph $G$, let $h(G)=\inf_{0<|A|<\infty}|\partial_G A|/|A|$, where $\partial_G A$ is the external vertex boundary. A $k$-net is a maximal set $N\subseteq V(G)$ whose distinct points have graph distance at least $k$. Its graph $G_k$ joins distinct points of $N$ whenever their original distance is at most $2k$.

For real $h>0$, integer $d\ge2$, and integer $k\ge1$, define the optimal uniform lower bound
\[
F(h,d,k)=\inf\{h(G_k):h(G)>h,\ \sup_v\deg_G(v)<d\},
\]
where the infimum ranges over all infinite connected graphs satisfying the conditions and all their $k$-nets. The fixed-$k$ comparison gives positive lower bounds. Is
\[
\inf_{k\ge1}F(1,10,k)>0?
\]
Equivalently, is there one $\varepsilon>0$ such that every net at every scale of every graph with $h(G)>1$ and maximum degree at most nine has Cheeger constant at least $\varepsilon$?

The source also asks for the behavior on a fixed regular tree $T_d$, $d\ge3$. With $F_d(k)=\inf_N h((T_d)_k)$ over its $k$-nets, is $\inf_kF_d(k)>0$, and does $F_d(k)\to\infty$ as $k\to\infty$? The infimum over nets is essential because maximal separated sets need not be unique or symmetric. These are uniform expansion questions; the degrees of the net graphs themselves may grow with $k$.
\subsection{Short English statement}
Do graph nets retain uniformly positive expansion at all scales, and does their expansion grow without bound on regular trees?
\subsection{Sources}
\begin{itemize}
\item[S1] ULAM.AI and the UnsolvedMath Contributors, supplied problems.json, record 10000012 (AMR-099-0012), AMR Open Problem Lists. Catalogue identity and source transcription; CC BY 4.0.
\url{https://huggingface.co/datasets/ulamai/UnsolvedMath}.
\item[S2] Benjamini - Coarse Geometry and Randomness (Saint-Flour notes), Open problem 1.72, PDF page 15. Original source indexed by AMR-099-0012.
\url{https://www.wisdom.weizmann.ac.il/~itai/stflouraug24.pdf#page=15}.
\item[S3] I. Benjamini, Coarse Geometry and Randomness, primary Saint-Flour notes; the numbered question and its surrounding definitions.
\url{https://arquivo.pt/noFrame/replay/20201231041548id_/http://www.wisdom.weizmann.ac.il/~itai/stflouraug24.pdf}.
\end{itemize}
\end{document}
